\section{Experiments}
\label{sec:experiments}

\textbf{All numbers in this section are preliminary, from a single v0.19
training run.} They are reported to show that the learned verdict head clears
the relevant baselines by a wide margin, not as final benchmark figures.

\subsection{Setup and metrics}

\paragraph{Evaluation set.}
We evaluate on a held-out set of $269$ trajectories: $130$ attacks and $139$
benign, drawn from multiple sources and frameworks.

\paragraph{Metrics.}
We frame metrics around deployment semantics rather than classification
(Section~\ref{sec:method:check}):
\begin{itemize}[leftmargin=1.4em,itemsep=0.1em]
  \item \textbf{Prevention Rate (PR):} on attack traces, the fraction flagged
        \emph{at or before} the anomaly step --- i.e. in time to actually block
        the harmful action. A flag after the anomaly step is a miss.
  \item \textbf{False-Stop Rate (FSR):} on benign traces, the fraction flagged
        at \emph{any} step.
  \item \textbf{PR@FSR$\le$1\%} (headline): the best achievable prevention rate
        subject to at most a $1\%$ false-stop rate, i.e. ``under one false stop
        per hundred benign runs, how many attacks do we prevent?'' We also
        report PR@FSR$\le$5\%.
  \item \textbf{AUROC} (reference) over trace-level scores, and
        \textbf{reason cosine similarity} between the decoded \texttt{<reason>}
        and the gold reason on correctly-stopped attacks.
\end{itemize}

\paragraph{Baselines.}
We compare against two intentionally simple scorers run through the same
harness: a \emph{keyword/regex} baseline (hand-written patterns on tool
arguments and observations) and an \emph{action-perplexity} baseline (raw
next-token perplexity over the proposed action span, using the base model
without fine-tuning). Both set a floor for ``did the learned head buy us
anything.''

\subsection{Main results}

\begin{table}[t]
  \centering
  \small
  \begin{tabular}{lccc}
    \toprule
    Method & PR@FSR$\le$1\% & PR@FSR$\le$5\% & AUROC \\
    \midrule
    Keyword / regex baseline      & $\approx 0.115$ & \TODO{n/a} & $\approx 0.56$ \\
    Action-perplexity baseline    & $\approx 0.08$  & \TODO{n/a} & \TODO{n/a} \\
    \textbf{IAM (verdict head)} & $\mathbf{\approx 0.89}$ & $\mathbf{\approx 0.95\text{--}1.0}$ & $\mathbf{\approx 0.997}$ \\
    \bottomrule
  \end{tabular}
  \caption{Preliminary v0.19 results on the 269-trace evaluation set
  ($130$ attacks / $139$ benign). The headline PR@FSR=1\% operating point uses
  a $P(\textsc{STOP})$ threshold $\tau\approx 0.92$ at the best epoch. Baselines
  are run through the identical harness. % TODO: fill baseline PR@FSR=5\% and action-perplexity AUROC.
  }
  \label{tab:main}
\end{table}

Table~\ref{tab:main} summarizes the main result. At an operating point of
$P(\textsc{STOP})$ threshold $\tau\approx 0.92$, IAM reaches
\textbf{PR@FSR=1\% $\approx 0.89$}; relaxing the constraint to FSR$\le$5\%
raises prevention to $\approx 0.95\text{--}1.0$. Trace-level separation is very
strong (\textbf{AUROC $\approx 0.997$}). Both baselines are far below: the
keyword baseline reaches PR@FSR=1\% $\approx 0.115$ (AUROC $\approx 0.56$, near
chance) and the action-perplexity baseline $\approx 0.08$. The learned,
densely-supervised verdict head thus improves prevention at the $1\%$ false-stop
operating point by roughly $8\times$ over the strongest baseline.

\paragraph{Per-position behavior.}
Decomposing by query position, the true-positive rate at the anomaly step
(attack-stop@anomaly) is $\approx 0.95\text{--}0.99$, while the in-trace
over-flag rate --- the model stopping at a \emph{safe} prefix step of an attack
trace (attack-safe) --- is only $\approx 0.02\text{--}0.03$. The dense
\textsc{OK} supervision is doing its job: the model stops at the right step
without flagging the benign lead-up.

\paragraph{Reason quality.}
On correctly-stopped attacks, the decoded \texttt{<reason>} attains
$\approx 0.92$ cosine similarity to the gold reason. We observe this degrades
slightly with more training epochs, consistent with mild overfitting on the
reason field.

\subsection{Per-family detection}
\label{sec:experiments:family}

Detection is essentially saturated ($\approx 1.0$) on the \emph{shaped} attack
families that dominate training --- exfiltration, email-injection, SSRF, shell
injection, persistence, and the sophisticated indirect-injection family
(\texttt{iif*}). It is weaker on families that are underrepresented or stylistically
distinct in our data, notably ATBench-native, MCP, and DNS-exfiltration
(\texttt{dnsx}) families. This pattern --- strong where supervision is dense,
weak where it is sparse --- is the clearest signal of where additional training
data is needed. % TODO: per-family detection table with counts; only qualitative summary available for v0.19.

\subsection{The bottleneck: ATBench benign false stops}
\label{sec:experiments:bottleneck}

The most important honest finding concerns \emph{where} the residual false
stops come from. The overall benign false-stop rate is $\approx 0.022$, but it
is \emph{not} spread across sources: it is concentrated almost entirely in a
single source. ATBench benign traces have an FSR of $\approx 0.065$, while
\emph{every other benign source has an FSR of $\approx 0.000$}. In other words,
the FSR floor --- and the binding constraint at the PR@FSR=1\% operating point
--- is a handful of ATBench benign traces that the model over-flags. Reducing
the false-stop rate further is therefore primarily a data problem on this one
source (more / cleaner ATBench-style benign coverage), not a change to the
method. We report this prominently because it, rather than attack recall, is the
current limiting factor for tightening the operating point.

\subsection{Red-team findings (context for positioning)}
\label{sec:experiments:redteam}

The following qualitative findings motivate the ``complement, not replacement''
positioning of Section~\ref{sec:intro}; they characterize \emph{when} a runtime
monitor is needed. % TODO: these are red-team probing observations from earlier (v0.2/v0.3) runs, reported qualitatively; not part of the 269-trace quantitative eval.

\begin{itemize}[leftmargin=1.4em,itemsep=0.15em]
  \item \textbf{F1 --- Aligned backends are robust.} A production-aligned model
        (Claude Code) resisted $0/65$ common agent attacks (indirect injection,
        exfiltration, SSRF, shell injection), with or without an added safety
        system prompt. A monitor is largely \emph{unnecessary} for such
        deployments.
  \item \textbf{F2 --- Sophistication tiers matter.} On a weakly-aligned
        backend, lazy indirect injection (HTML comment, dev note) had a
        $\approx 0\%$ bite rate, but sophisticated patterns
        (false-completion, multi-step protocol, persona override) reached
        $\approx 78\%$. A single ``indirect-injection rate'' is misleading;
        benchmarks should stratify by tier.
  \item \textbf{F3 --- Temperature is an attacker dimension.} Multi-step
        injection bite rate rose from $\approx 0\%$ at temperature $0$ to
        $\approx 50\%$ at $0.3$; single-step attacks (SSRF, direct exfil) were
        temperature-invariant. Greedy-decoding red teams under-report real
        feasibility.
  \item \textbf{F4 --- LLM-designed attacks beat hand-authored.}
        LLM-synthesized scenarios (e.g. sandbox-escape, output-manipulation)
        reached $\approx 83\text{--}100\%$ success in several families, above
        the $\approx 50\%$ of comparable hand-authored scenarios.
  \item \textbf{F5 --- CVE patterns replicate.} A reproduction of a zero-click
        email-injection pattern (CVE-2025-32711) had a $\approx 100\%$ bite rate
        on a weakly-aligned backend; a PR-description RCE pattern
        (CVE-2025-53773) reached $\approx 44\%$.
  \item \textbf{F6 --- Prompt extraction is hard.} Direct attempts to leak a
        sentinel-laden system prompt failed ($0/25$); even weakly-aligned models
        have specific coverage for system-prompt protection.
  \item \textbf{F7 --- Canonicalization is feasible.} The same scenarios run
        across distinct frameworks fold into one canonical schema, and patterns
        defined over it transfer (Section~\ref{sec:data:canon}).
  \item \textbf{F8 --- Attack success is 2D.} Attack success is governed jointly
        by attacker goal and victim alignment: near $0\%$ across goals for a
        high-alignment backend, but $60\text{--}96\%$ for a medium-alignment one.
        IAM targets the weakly-aligned cell of this matrix.
\end{itemize}

\noindent
\TODO{latency benchmarks (p50/p99 per check on GPU and CPU) --- not yet
measured.}\quad
\TODO{quantitative cross-backend transfer (train-on-A / test-on-B) --- not yet
measured.}\quad
\TODO{ablations: with/without dense supervision, with/without $W_{\textsc{STOP}}$
up-weighting, with/without anomaly truncation --- not yet run.}
